\documentclass[12px]{article}

\title{Lezione 1 Metodi Numerici di Approssimazione}
\date{2026-02-27}
\author{Federico De Sisti}

\input{../../setup.tex}

\begin{document}
	\maketitle
	\newpage
	\subsection{Prima parte della lezione (da recuperare)}
	\subsection{Seconda parte della lezione}
	se $\frac km := \omega^2$ otteniamo
	 \[
	\ddot y + \omega ^2 y = 0
	.\] 
	Otteniamo l'equazione dell'oscillatore armonico.\\
	Possiamo per esempio vedere una corda di chitarra come una serie di oscillatori armonici.
\subsection{Da sistema non autonomo a sistema autonomo}
\[
\begin{cases}
	\dot y(t) = f(t,y(t))\\
	y(0) = y_0
\end{cases}
\]
Questo sistema può essere trasformato nella forma
\[
\dot Y = F(Y)
.\] 
Proviamo a pensare che $t$ dipenda da un tempo assoluto. Vogliamo scrivere un ODE per il tempo, come varia il tempo rispetto ad un altra variabile ausiliaria.\\
Chiamiamo $t = t(s)$ tale variabile, e voglio descrivere come varia $t$ in base ad $s$.

Per esempio  $\dot t(s) = 1$, ovvero percepiamo il tempo con una pendenza costante verso il futuro. 

Questa si risolve facilmente con $t(s) = s + const$,  $t(0) = 0 \Rightarrow const = 0 \Rightarrow  t(s) = s$ \\

Combino questa variabile ausiliaria nel sistema precedente
\[
\dot y(t(s)) = f(t(s),y(t(s)))
\]\[\dot t (s) = 1
.\] 
$Y(s) = (t(s), y(s))$ \\
$\dot Y(s) = (\dot t(s), \dot y(s)) = (1,f(t(s),y(t(s)))) = F(Y)$ \\

Così posso eliminare la dipendenza dal tempo dal sistema ed avere solamente la dipendenza dalla variabile $Y$.\\

Possiamo quindi trattare solo sistemi autonomi in quanto possiamo ricondurci.
\newpage
\begin{defi}[Lipschitziana]
	$f$ si dice \textbf{lipschitziana} in $(t_0,y_0)$ rispetto ad $y$ se $\exists L > 0 $ tale che
	\[
	|f(t,y_1)-f(t,y_2)|\leq L |y_1-y_2|
	.\] 
	$\forall t\in B_\varepsilon(t_0)$\\
 $\forall y_1,y_2\in B_\delta(y_0)$ intorni
\end{defi}
Se $y_1\neq y_2$ 
\[
	\left|\frac{f(t,y_1) - f(t,y_2)}{y_1-y_2}\right| \xrightarrow
{y_1 \rightarrow y_2 \rightarrow \bar y} \left|\frac{\partial f(t,y)}{\partial y}\right|\leq L.\] 
Esempi di funzioni Lipschitziana sono $x$, $x^2$ ma solamente in un compatto, $arctan(x)$. 
\begin{defi}[ODE]
	Definiamo $ODE$ il seguente sistema
	\[
	\begin{cases}
		\dot y(t) = f(t,y(t)) \ \ t > 0\\
		y(0) = y_0
	\end{cases}
	.\] 
\end{defi}
\begin{teo}[Teorema di esistenza e unicità di una soluzione locale di un problema di Cauchy]
	Se $f$ è localmente lipschitziana $ \Rightarrow $ $\exists !$ soluzione locale di  $ODE$
\end{teo}
\textbf{Equazione di Volterra}\\
La soluzione che cerchiamo possiamo scriverla come 
\[
y(t) = y_0 + \int^t_0f(\tau,y(\tau))d\tau
.\] 
Se $f$ sia continua e $y$ è continua, la composizione è continua, in particolare la funzione è integrale è  $C^1$, quindi posso applicare il teorema fondamentale del calcolo integrale.
 \[
\dot y(t) = F'(t) = f(t,y(t))
.\] 
Solitamente si riferisce a questa forma come forma debole della soluzione.\\

Le due formulazioni sono le due strade che prenderemo per trovare schemi numerici che risolvono i sistemi.\\

Possiamo quindi migliorare la qualità della soluzione migliorando la qualità dell'approssimazione della derivata o dell'integrale.

\subsection{Soluzioni globali definite per $t\in (0,+\infty)$}
Con $f$ Lipschitziana globale, avremmo soluzione globale, ma poche funzioni sono di questa classe.\\

Esempio $\dot y = y(1-y)$, non è globalmente lipschitziana poiché il membro sulla destra è quadratico, Esempio dello sciacquone, più acqua è presente, meno entra successivamente fino a fermarsi. \\


\end{document}
